\section{Inertia.js}
Komunikace mezi serverovou částí aplikace a klientským rozhraním je zprostředkována knihovnou Inertia.js. Jedna ze základních funkcionalit je poskytována metodou \texttt{Inertia::render}, která se v controllerech využívá pro vytvoření speciální Inertia odpovědi. Její využití je demonstrováno na výpisu kódu \ref{listing:inertia-example}. Prvním z parametrů této metody je název Vue.js komponenty, která má být vykreslena. Jako druhý parametr je očekáváno asociativní pole dat, která jsou do komponenty předána jako její vlastnosti.

Vzhledem k vyšší komplexitě databázových operací byla při implementaci controlleru pro hlavní dashboard klíčová minimalizace počtu prováděných SQL dotazů. Klíčovým prvkem tohoto řešení je využití anonymních funkcí u hodnot jednotlivých vlastností. Na výpisu kódu \ref{listing:inertia-example} je tento přístup využit u klíčů \texttt{users} a \texttt{developerStats.data}. Díky tomuto zápisu nedochází k okamžitému provedení dotazů při každém volání controlleru, ale je využito mechanismu odloženého vyhodnocení. Ten je nezbytný pro efektivní fungování částečného vykreslování.

V praxi to znamená, že pokud uživatel například změní počet záznamů na~stránce u jedné z tabulek, tak Inertia.js odešle požadavek obsahující hlavičku \texttt{X-Inertia-Partial-Data}, která obsahuje unikátní identifikátor dané tabulky. Server následně na základě této hlavičky vrátí pouze ta data, která jsou potřebná pro její aktualizaci. Díky odloženému vyhodnocení nejsou ostatní databázové dotazy v rámci tohoto požadavku vůbec spuštěny, což výrazně zrychluje odezvu aplikace.

\begin{listing}[h!]
    \caption{Předávání dat knihovně Inertia.js}\label{listing:inertia-example}
    \begin{minted}{php}
<?php

private const string TABLE_DEVELOPER_STATS = 'developerStats';

// ...

return Inertia::render('Dashboard', [
    'users' => fn () => $usersQuery->get()->toArray(),
    'from' => $from->toDateString(),
    'to' => $to->toDateString(),
    self::TABLE_DEVELOPER_STATS => [
        'data' => fn () => $this->getDeveloperTableStats(
            $request,
            $vcsInstanceUserIds,
            $from,
            $to
        ),
        'config' => [
            'id' => self::TABLE_DEVELOPER_STATS,
            'pageParamName' => self::PAGE_PARAM_NAME,
            'perPageParamName' => self::PER_PAGE_PARAM_NAME,
        ],
    ],
    // ...
]);
    \end{minted}
\end{listing}
